\documentclass[11pt,a4paper]{article}
\usepackage[margin=25mm]{geometry}
\usepackage{amsmath,amssymb,amsthm,booktabs,hyperref}
\setlength{\emergencystretch}{2em}
\newtheorem{lemma}{Lemma}\newtheorem{theorem}[lemma]{Theorem}
\newtheorem{proposition}[lemma]{Proposition}
\newcommand{\gen}[1]{\langle #1\rangle}
\newcommand{\PSL}{\operatorname{PSL}}
\title{Absolute flags and their moving lines}
\author{Human Compression project: L017 working recovery}\date{}
\begin{document}\maketitle
This author manuscript proposes a corrected recovery of results of Dimitri Leemans, Klara Stokes
and Philippe Tranchida, \emph{On trialities and their absolute geometries},
21 March 2024. The geometries and claimed parameters belong to those authors.
The coordinate argument below replaces their isolated small twisted-field
computation. Its complete source contract and finite certificates are supplied
separately in the same packet. Distinct review and whole-original acceptance
remain pending. No numerical factor or formalization is claimed.

\section{Classical input and conventions}
Let $F$ be a finite field, $\sigma$ an automorphism of order one or three,
$K=\operatorname{Fix}(\sigma)$, $f=|F|$ and $k=|K|$. In $\mathrm{PG}(7,F)$ put
\[
 Q(X)=X_0X_4+X_1X_5+X_2X_6+X_3X_7.
\]
The $D_4$ geometry has singular points, the two families of maximal singular
three-spaces, and singular lines. Opposite-family three-spaces are incident
when their intersection is a plane. All other incidence is containment.

We retain the following precise classical input from Tits's triality theory,
as cited in source Section 2.3 and reference [17], and the generalized-polygon
facts cited there from [18]. A triality of type $I_\sigma$ is conjugate to the
canonical cyclic triality followed by coordinatewise $\sigma$. Its fixed
geometry is a Moufang generalized hexagon of order $(f,k)$: each absolute
line has $f+1$ points and each absolute point lies on $k+1$ absolute lines.
The commuting projectivity group $G$ is $G_2(k)$ or ${}^3D_4(k)$ respectively.
It is transitive on apartments with a distinguished chamber, and on ordered
point paths $(A,P,B)$ of length two in the point-collinearity graph, with
the two lines distinct. Any two points lie in an apartment.
These classification, embedding and transitivity assertions are retained
inputs, not newly proved Lie-theoretic results.

The incidence graph of this hexagon has girth twelve and diameter six.
An apartment therefore has six points and six lines. When we say that two
points are at distance $j$ in the hexagon, we mean point-collinearity distance;
their incidence-graph distance is $2j$. All distances in the moving geometry
will instead be incidence-graph distances. This distinguishes the two
conventions used in the source's counting proof.

Here is the explicit canonical coordinate model needed from the classical
input. In three copies of the point coordinates set
\begin{align*}
T(X,Y,Z)={}&\det\begin{pmatrix}X_0&X_1&X_2\\Y_0&Y_1&Y_2\\Z_0&Z_1&Z_2\end{pmatrix}
+\det\begin{pmatrix}X_4&X_5&X_6\\Y_4&Y_5&Y_6\\Z_4&Z_5&Z_6\end{pmatrix}\\
&+X_3(Z_0Y_4+Z_1Y_5+Z_2Y_6)+X_7(Y_0Z_4+Y_1Z_5+Y_2Z_6)\\
&+Y_3(X_0Z_4+X_1Z_5+X_2Z_6)+Y_7(Z_0X_4+Z_1X_5+Z_2X_6)\\
&+Z_3(Y_0X_4+Y_1X_5+Y_2X_6)+Z_7(X_0Y_4+X_1Y_5+X_2Y_6)\\
&-X_3Y_3Z_3-X_7Y_7Z_7.
\end{align*}
A point $X$ is incident with the type-one object represented by $Y$ when
$T(X,Y,Z)=0$ for every $Z$; the other incidences follow cyclically.
The canonical triality cycles the three coordinate copies. Its semilinear
version applies $\sigma$ as well. In particular an absolute point satisfies
$T(X,X^\sigma,Z)=0$ for every $Z$. Applying the triality supplies the other
two required incidences.

\section{Special planes and a uniform distant pair}
Write $e_i$ for the coordinate points. At $P=e_0$, direct comparison of
coefficients in the two incidence equations gives
\[
 P^\alpha=\gen{e_0,e_3,e_5,e_6},\qquad
 P^{\alpha^2}=\gen{e_0,e_5,e_6,e_7},\qquad
 P^\omega=\gen{e_0,e_5,e_6}.
\]
For $X=a e_0+b e_5+c e_6$ its absolute-point condition is precisely
$bc^\sigma-cb^\sigma=0$. Thus the absolute points of this plane are the
union of the $k+1$ pencil lines through $P$ with $[b:c]\in\mathrm{PG}(1,K)$.
Each pencil line has all its $f+1$ points absolute. These are exactly the
absolute lines through $P$, by the retained valency. Point transitivity
transfers this description to every $P^\omega$.

The following coordinate apartment is defined over $K$:
\[
                  e_0,e_5,e_2,e_4,e_1,e_6,e_0.             \tag{1}
\]
For completeness, in the untwisted hyperplane $X_7=-X_3$ the fixed lines
are characterized by the source's six Pl\"ucker equations
\[
X_{12}=X_{34},\quad X_{54}=X_{32},\quad X_{20}=X_{35},\quad
X_{65}=X_{30},\quad X_{01}=X_{36},\quad X_{46}=X_{31}.       \tag{2}
\]
The six lines in (1) satisfy (2), so are fixed by the canonical triality;
their $K$-definition makes them fixed by the semilinear triality too.
The opposite pairs $(e_0,e_4),(e_5,e_1),(e_2,e_6)$ have nonzero polar pairing
for $Q$. Apartment transitivity now proves that opposite absolute points
are never joined by a singular line. This proves source Lemma 3.1 without
its ambiguous subspace argument.

\begin{lemma}\label{plane}
Every moving line contains exactly $k+1$ absolute points and lies in a
unique special plane $P^\omega$, without passing through $P$.
Conversely the $f^2$ lines of $P^\omega$ not through $P$ are moving lines.
\end{lemma}
\begin{proof}
Take two absolute points on a moving line and put them in an apartment.
They are not adjacent, since their projective join would then be fixed;
they are not opposite by the preceding polar calculation. They therefore
have a unique middle point $P$ on a two-edge point path. Their join lies
in $P^\omega$ and avoids $P$. It meets each of the $k+1$ absolute pencil
lines once, in distinct points, and these are all its absolute points by
the coordinate description. Conversely any nonpencil line has those
$k+1$ absolute points and cannot be fixed: otherwise it and two pencil
lines would give a triangle in the generalized hexagon.
If it lay also in $R^\omega$, any two of its absolute points would have
both $P$ and $R$ as middle points. Girth twelve implies $P=R$.
A projective plane over $F$ has exactly $f^2$ nonpencil lines.
\end{proof}

\begin{lemma}\label{far}
The two moving lines $l=\gen{e_5,e_6}$ and $l'=\gen{e_1,e_2}$
have incidence-graph distance at least six, for both $f=k$ and $f=k^3$.
\end{lemma}
\begin{proof}
They are nonpencil lines in $e_0^\omega$ and $e_4^\omega$.
Their absolute points are exactly their $K$-rational points: these already
give $k+1$ absolute points, the full number from Lemma~\ref{plane}.
Let $A=a e_5+b e_6$ and $B=c e_1+d e_2$ be such points. Their join is
singular exactly when $ac+bd=0$. All six equations (2) for this join
are $0=0$. Thus, if singular, it is fixed by the canonical triality;
being defined over $K$, it is also fixed by the semilinear triality.
Consequently no moving line joins an absolute point of $l$ to one of $l'$.
The two original lines are disjoint, so there is no path of length zero,
two or four between them. This includes $(k,f)=(2,8)$ without a finite
computation.
\end{proof}

The unnumbered source assertion about the image of a moving line also
has a direct proof. For $l=\gen{e_5,e_6}$ the incidence equations give
$\alpha(l)=\gen{e_0,e_7}$. Substituting $X=a e_0+b e_7$ in
$T(X,X^\sigma,Z)$ makes the coefficient of $Z_7$ equal to $-bb^\sigma$,
so $e_0$ is its only absolute point. The transitivity proved next, and
the fact that $G$ commutes with $\alpha$, transfer this to every moving
line. In particular $\alpha$ does not act on the chosen moving-line set.

\section{All parameters and distances}
The retained transitivity on ordered two-edge point paths maps any moving
flag $(A,l)$ to any other: choose $B\ne A$ on $l$ and let $P$ be their
middle point. The path $(A,P,B)$ determines the flag. It also shows that
$G_P$ is transitive on all $f^2$ nonpencil lines in $P^\omega$, recovering
source Lemma 3.12.

The number of points at point distance two from an absolute point is
$k(k+1)f^2$, by the unique paths of length four in the hexagon. Each moving
line through that point accounts for exactly $k$ of them, so its moving
valency is
\[
                              C=(k+1)f^2.                 \tag{3}
\]
Counting successive levels of a generalized hexagon of order $(f,k)$
gives
\[
             N=(f+1)(1+fk+f^2k^2)                         \tag{4}
\]
points. One can check the last level directly: the first two nonzero point
levels have sizes $f(k+1)$ and $f^2k(k+1)$; the last has $f^3k^2$, since
each last-level point is reached through all its $k+1$ lines. Summing gives
(4). Double-counting flags with (3) yields $Nf^2$ moving lines.
The elementary breadth-first bound from a moving line is exactly
\[
               1+(k+1)(C-1)+k(k+1)(C-1)^2                \tag{5}
\]
lines within distance four. Repeated vertices can only lower this bound.
This recovers source Lemma 3.8, although Lemma~\ref{far} makes its
comparison with the total unnecessary for the diameter proof.

\begin{theorem}\label{classical}
The moving geometry is flag-transitive, has $N$ points and $Nf^2$ lines,
line size $k+1$ and point valency $(k+1)f^2$. Its incidence girth is six,
point eccentricity is five, and line eccentricity is six.
\end{theorem}
\begin{proof}
The geometry has no isolated elements by Lemma~\ref{plane} and (3).
Distinct projective lines have at most one point in common, so its incidence
girth is at least six. Alternate points of the apartment (1) form a
noncollinear triangle of moving joins, giving equality.

Every two absolute points have moving incidence distance at most four.
For point distance two in the hexagon their moving join suffices.
For adjacent points $A,B$, choose a third point $P$ on their absolute line
and an absolute point $D$ on another absolute line through $P$.
Both $DA,DB$ are nonpencil lines of $P^\omega$, hence moving.
For opposite points, label an apartment cyclically by $A,P_1,P_2,B,P_4,P_5$.
Choose a third point $D$ on the line $P_1P_2$. Then $AD$ lies in
$P_1^\omega$ and $BD$ in $P_2^\omega$; both are nonpencil, hence moving.

Every point is consequently within five of every moving line, and all pairs
of moving lines are within six. The latter bound is attained by
Lemma~\ref{far}. To attain the former, take a point $A$ of an apartment
and the moving join of its neighbor $P_1$ with its opposite point $B$.
The quadric's all-or-one property, or the polar linear equation on this
line, says its only point singular-collinear with $A$ is $P_1$, because
$A$ is not singular-collinear with $B$. But $AP_1$ is fixed. There is
therefore no moving path of length three from $A$ to that moving line.
The distance is five. Point and line transitivity give the asserted
eccentricities at every element.
\end{proof}

The subgroup implementation following source Lemma 3.12 requires an actual
geometric identification. A stabilizer of a moving line in $P^\omega$ has
index $f^2$ in $G_P$, but an abstract subgroup of that index is not thereby
identified with the stabilizer. A correct implementation constructs the
line and its stabilizer, or verifies its orbit and incidence in the
coordinate geometry. Orders and indices alone are insufficient.

\section{Completing the even projective-linear construction}
This section supplies the characteristic-two counterpart of the paid
odd-characteristic signed-vector theorem in the accepted L023 recovery.
It uses no finite computation.

Let $F$ have characteristic two and let $a,c,b\in\mathrm{SL}_2(F)$ be
involutions, with $a,c$ distinct and commuting, such that $b$ does not
commute with any of $a,c,ac$. Suppose $\gen{a,b,c}=\PSL_2(F)$.
Distinct commuting nonidentity unipotents have the same fixed point on
$\mathrm{PG}(1,F)$; this follows directly from the centralizer of
$\left(\begin{smallmatrix}1&x\\0&1\end{smallmatrix}\right)$.
Conjugate their fixed point to infinity. The matrix $b$ cannot fix
infinity, since then the generated group would be triangular. An upper
unitriangular conjugation makes its diagonal entries zero, and a diagonal
conjugation makes its lower-left entry one, using the unique square root
in $F$. We therefore have
\[
 a=u(x),\quad c=u(y),\quad ac=u(x+y),\quad
 u(z)=\begin{pmatrix}1&z\\0&1\end{pmatrix},\qquad
 b=J=\begin{pmatrix}0&1\\1&0\end{pmatrix},                 \tag{6}
\]
where $x,y,x+y$ are nonzero and distinct.
Put $x_i=x,x_j=y,x_k=x+y$ and $H_i=\gen{u(x_i),J}$, and similarly
for $j,k$. Both generators of $H_i$ preserve the normalized quadratic form
\[
 F_i(s,t)=a_i(s^2+t^2)+st,\qquad a_i=x_i^{-1}=t_i^2.
\]
All normalized binary quadratic forms have unique coordinates
$F_{u,v}=u^2s^2+st+v^2t^2$. The action of $\mathrm{SL}_2(F)$ on them
is faithful: preserving all their differences forces the action on every
squared linear form to be trivial, hence the original linear action is
trivial.

\begin{lemma}\label{evenpair}
The simultaneous stabilizer of $F_i,F_j$ is $\{1,J\}$. Consequently
$H_i\cap H_j=H_i\cap H_k=H_j\cap H_k=\gen J$.
\end{lemma}
\begin{proof}
Their difference is $(a_i+a_j)(s+t)^2$, with nonzero coefficient.
A simultaneous stabilizer therefore fixes the linear form $s+t$.
A determinant-one matrix doing so has the form
$\left(\begin{smallmatrix}p&1+p\\1+p&p\end{smallmatrix}\right)$.
Preserving $F_i$ imposes $p^2+p=0$, so it is $I$ or $J$.
These two matrices lie in both parabolics; the same proof works for every pair.
\end{proof}

\begin{lemma}\label{evenconics}
The orbit of $F_j$ under $H_i$ and its orbit under $H_k$ intersect
in at most three quadratic forms.
\end{lemma}
\begin{proof}
For an isometry of $F_i$, the dual action on linear forms preserves
the dual form. Its coefficients are the same as those of $F_i$, since
the polar matrix interchanges the two coordinates. Also
$F_j-F_i=(t_j+t_i)^2(s+t)^2$, and $F_i(1,1)=1$.
It follows that any $F_{u,v}$ in the first orbit satisfies
\[
 a_i(u+v)^2+uv+t_i(u+v)=a_j.                              \tag{7}
\]
Explicitly apply $F_i$ to $(u+t_i,v+t_i)$ and set its value equal
to $(t_j+t_i)^2$; expansion gives (7). The second orbit satisfies
the same equation with $i$ replaced by $k$. Subtraction, with $w=u+v$,
gives
\[
 (t_i+t_k)^2w^2+(t_i+t_k)w=0.
\]
Since $t_i\ne t_k$, either $w=0$ or $w=(t_i+t_k)^{-1}$.
The first choice gives $u=v=t_j$ uniquely. In the second case, substitute
$v=u+w$ in (7); the resulting polynomial in $u$ has degree two and
leading coefficient one. It has at most two roots. We need only this
upper bound, not a supposition that both roots are present.
\end{proof}

Consider the coset pregeometry with parabolics $H_i,V,H_j,H_k$, where
$V=\gen{a,c}$. Their intersections with $V$ are respectively
$\gen a,\gen c,\gen{ac}$: for example $u(y)$ preserves $F_i$ only if
$a_i y^2+y=0$, namely $y=0$ or $y=x_i$.
In the residue of $H_i,H_k$, a vertex coset $gH_j$ is sent to
the form $F_j^{g^{-1}}$, with action convention chosen so that left
cosets give a well-defined form. Represent $g$ in $H_i$ to see injectivity
on this face-star from Lemma~\ref{evenpair}. Represent it also in $H_k$
to apply Lemma~\ref{evenconics}. Thus there are at most three vertices.
The three cosets
\[
                        H_j,\quad aH_j,\quad baH_j
\]
are present and distinct, just as in the odd signed-vector argument:
the necessary intersections use $ac$ and $bac$, and equality of the last
two would force $aba=b$. There are exactly two incident edges, $V,bV$.
Indeed $V=\gen a\gen{ac}$ permits its two factors to be absorbed into
representatives in $H_i,H_k$; their intersection is $\{1,b\}$.
These edges join the first vertex to each of the other two.

\begin{theorem}\label{evencompletion}
The characteristic-two pregeometry is a geometry. Each face--Petrie
residue is the path with three vertices and two edges; it has two free
group orbits of chambers and is not firm.
\end{theorem}
\begin{proof}
The preceding calculation completes every flag of the three nonedge types.
Every flag containing an edge completes because its three other types
form a complete three-partite structure, with two choices in each part.
Incident pairs avoiding edges have common representatives and also extend.
These are all flags. At an edge the eight chambers split into two free
$V$-orbits, and hence globally into two free group orbits. The two endpoints
of the displayed residue each have just one completing edge, proving
nonfirmness. The exact pair intersections also prove the usual faithful
thin rank-three map conditions; they are not inferred from an arbitrary
quotient of the universal map group.
\end{proof}

\section{The fixed absolute of a cubic projective-linear marking}
Let $G=\PSL_2(q^3)$ and let $a,b,c$ satisfy the five hypotheses of source
Section 3.2: they are generating involutions, $a,c$ commute, an actual
order-three outer automorphism $\alpha$ cycles $a,c,ac$ and fixes $b$,
and no marked duality exists. Writing $\alpha\in\operatorname{Out}(G)$
in the source must mean an actual automorphism representing that outer
class; otherwise its action on elements and cosets is undefined.
The three periods $|ab|,|cb|,|(ac)b|$ are equal and exceed two.
If they were at most two, all three generators would commute, contrary
to generation of this nonabelian simple group.

The odd-characteristic completion theorem retained in full from the
accepted L023 packet, and Theorem~\ref{evencompletion} in characteristic
two, apply. In particular all flags complete, all three outer-parabolic
pair intersections are $\gen b$, and the exact pair residue is a path
with three vertices and two edges. The rank-three truncation is the faithful
map obtained by gluing its flag triangles: all its panels are thin, its
flag-adjacency graph is connected by generation, and its vertex and face
residues are the marked dihedral cycles. Unique flag adjacencies imply that
a map automorphism fixing a flag fixes every flag. The regular $G$-action
is therefore its full map automorphism group.

This also justifies the source's triality-without-duality assertion for
these particular four-type geometries. The edge type has $|G|/4$ objects,
whereas every other type has $|G|/(2m)$, with $m>2$ the common period.
Any correlation fixes the edge type. If it induced a transposition on the
remaining three types, conjugating by the known triality would give a
vertex--face interchange fixing the Petrie type. Its restriction to the
faithful map would be a duality. Compose with an element of the flag-regular
group to place its image base flag back at the base flag; the induced marked
automorphism of $G$ then interchanges $a,c$ and fixes $b$, contrary to source
hypothesis (5). Thus the exact type-permutation image is $C_3$.

We retain the standard projective-linear
descent fact used in the source and the accepted cubic-generation proof:
an order-three outer automorphism of $\PSL_2(q^3)$ is conjugate in
$\operatorname{PGL}_2(q^3)$ to one of the two generating field automorphisms,
and its fixed group is isomorphic to $\PSL_2(q)$. This uses the odd degree
of the field extension and is not a consequence merely of the word ``outer.''

\begin{theorem}\label{fixed}
The fixed absolute, whose objects are minimal invariant flags for the two
type orbits, is a disjoint union of $|\PSL_2(q)|/2$ paths of length two.
It has $|\PSL_2(q)|$ fixed edges and $3|\PSL_2(q)|/2$ absolute outer objects.
\end{theorem}
\begin{proof}
Use left cosets, with $H_i=\gen{a,b}$, $H_j=\gen{c,b}$,
$H_k=\gen{ac,b}$ and $V=\gen{a,c}$. Every stable coset of $V$ has exactly
one $\alpha$-fixed element. Indeed its four elements split into orbits
of lengths one or three. Four fixed elements would make all their quotients
in $V$ fixed, whereas $\alpha$ cycles its three nonidentity elements.
Conversely a fixed element gives a stable coset. Thus fixed edges are
in bijection with $C_G(\alpha)$.

At the fixed edge $V$, the triality cycles three two-element outer-type
sets, so there are exactly two invariant outer triples. One is
$(H_i,H_j,H_k)$, whose common edges are $V,bV$ by the exact pair residue.
Call such a triple a center. The other is
\[
                         (cH_i,acH_j,aH_k).               \tag{8}
\]
It has no common group representative. In fact the intersection of the
first two cosets is $\{ac,acb\}$, and neither lies in $aH_k$, since
$c\notin H_k$ and $b\in H_k$. Call such a triple a leaf.

Every outer triple extends to an edge and is a translate of one of these
two triples, because the eight local chambers form two free $V$-orbits.
These two classes cannot coincide: existence of a common representative
is invariant under translation. The center has two completing edges.
The leaf has just one. To verify the latter precisely, translate (8) to
$(H_i,H_j,cH_k)$. Its stabilizer is contained in $H_i\cap H_j=\{1,b\}$.
If $b$ stabilized the last coset, then $cbc$ would lie in $H_k\cap H_j$
and equal $b$, contrary to $|bc|>2$. Its stabilizer is therefore trivial.
There are $|G|$ leaf triples and $|G|$ chambers in their chamber orbit;
each leaf consequently has one completing edge.

Any invariant outer triple has one or two completing edges. The order-three
automorphism fixes each of them, so there are no additional isolated
absolute objects. Every fixed edge meets one center and one leaf, as at
$V$; its fixed representative translates the local calculation by an
element commuting with $\alpha$. Centers have degree two and leaves
degree one. The absolute is therefore a disjoint union of two-edge paths.
Its number of edges is $|C_G(\alpha)|=|\PSL_2(q)|$, proving the formula.
\end{proof}

This proof also gives a small exact construction for the moving absolute.
Let $H=C_G(\alpha)$. All fixed edges are $hV$ with $h\in H$, and every
absolute outer object is incident with one of them. Its complete candidate
list is therefore the two outer cosets at each $hV$, with duplicates removed.
For each such object, enumerate its incident original edges and keep an
edge precisely when it is moved by $\alpha$ and its other endpoint is also
absolute. This is an exhaustive construction, not a heuristic sample of
the large group. Isolated vertices can then be removed using the source's
stated convention.

The symmetric-group centralizer in source Step 2 cannot supply the desired
triality by conjugation: every element of that centralizer induces the
identity automorphism of $G$. Instead use the explicitly given semilinear
field automorphism (or the normalizer with the action verified). Also
$\operatorname{Aut}(\PSL_2(q^3))$ is not uniformly an extension of order
three of $G$: the full field and diagonal automorphisms must be accounted
for. The incidence and edge-fixedness tests in Steps 3--5 remain valid
once this actual automorphism and the correct parabolics are supplied.

\end{document}
